% Rules set straight into a paragraph, not in a box: a bare \vrule's width is part of its line, and a tall one opens the line spacing as TeX opens it.
\documentclass{article}
\begin{document}
A paragraph whose own list holds rules, as a plain \verb|\vrule| in running
text leaves them: a bar three points wide \vrule width 3pt height 1ex\relax{}
here, another one \vrule width 5pt height 0.5ex depth 0.5ex\relax{} there, and
a third \vrule width 2pt\relax{} that runs to the height of its line. Each
takes its width on the line, so the words after it move along by as much,
and the line breaks where TeX breaks it, at every width.

A tall rule \vrule width 0.4pt height 18pt depth 7pt\relax{} reaches above
and below its line, so TeX sets that line further from the one before and
the one after it than the others, which keep their baseline skip. The
words around it go on long enough for the rule to land on another line at
each width the test tries, and a second one \vrule width 1pt height 15pt
depth 1pt\relax{} later on does the same, a few lines further down.
\end{document}
